\chapter{Нормировка трёхчастичного шумового подъёма}
\label{ch:v8-baryon-three-particle-lift-normalization-gate}

Одночастичный направленный генератор фиксирует оператор потерь на $Q_L$,
но из этого ещё не следует единственный способ коррелировать три его копии.
Проверим данную свободу точно, не вводя чисел с плавающей точкой.

\section{Перестановочно-инвариантная ковариация копий}

Для каждого одночастичного скачка разрешим общую ковариацию трёх копий
\begin{equation}
 R_c=(1-c)I_3+c\mathbf1\mathbf1^{\mathsf T}
 =\begin{pmatrix}1&c&c\\c&1&c\\c&c&1\end{pmatrix}.
 \label{eq:v8-baryon-copy-covariance}
\end{equation}
Её собственные значения равны
\begin{equation}
 1-c,\quad1-c,\quad1+2c,
 \label{eq:v8-baryon-copy-covariance-spectrum}
\end{equation}
поэтому полная положительность эквивалентна
\begin{equation}
 -\frac12\leq c\leq1.
 \label{eq:v8-baryon-copy-correlation-cone}
\end{equation}
Диагональ $R_c$ тождественно единична. Следовательно, все значения $c$ в
этом отрезке дают один и тот же генератор на одночастичных наблюдаемых.
Точка $c=0$ описывает независимые среды, а $c=1$ --- общий коллективный шум.

\begin{theorem}[Неединственность трёхчастичного подъёма]
Одночастичный КМС-генератор вместе с перестановочной симметрией трёх копий
не выбирает трёхчастичный шумовой оператор. Остаётся непрерывный вполне
положительный отрезок $c\in[-1/2,1]$ с одинаковой одночастичной динамикой.
\end{theorem}

\section{Происхождение коэффициентов три и три вторых}

На слабом пространстве $(\mathbb C^2)^{\otimes3}$ положим
\begin{equation}
 \mathbf I=\sum_{r=1}^{3}\mathbf t^{(r)},
 \qquad
 \mathbf t=\frac12(\sigma_1,\sigma_2,\sigma_3).
 \label{eq:v8-baryon-total-isospin}
\end{equation}
Тогда
\begin{equation}
 \mathbf I^2=\frac{15}{4}P_{\rm sym}+\frac34P_{\rm mix},
 \qquad
 \rank P_{\rm sym}=\rank P_{\rm mix}=4.
 \label{eq:v8-baryon-isospin-projectors}
\end{equation}
Аддитивная часть содержит три одинаковых одночастичных вклада и потому
равна $K_3=3J_{1q}$; это точный тензорный подсчёт. Разность слабых уровней
для корреляции $c$ равна
\begin{equation}
 \Delta_3(c)=3c\kappa_2,
 \qquad
 \kappa_2=\frac{20x+22}{x+5}.
 \label{eq:v8-baryon-three-copy-splitting}
\end{equation}
Для двух копий аналогично $\Delta_2(c)=2c\kappa_2$, поэтому
\begin{equation}
 \frac{\Delta_3(c)}{\Delta_2(c)}=\frac32
 \qquad(c\ne0).
 \label{eq:v8-baryon-three-over-two-ratio}
\end{equation}
Таким образом, коэффициенты $3$ и $3/2$ действительно структурны, но
выражают соответственно число копий и разность казимиров слабого изоспина.
Они сами по себе не выбирают физическую корреляцию $c$.

\section{Точное семейство дискриминаторов}

После исправления одночастичной переносной нормы предыдущей главы
получаем
\begin{equation}
 v_c(x)=c\,
 \frac{52(25x^2+38x+13)}
 {375x^3+3916x^2+7267x+2782}.
 \label{eq:v8-baryon-correlation-family-discriminator}
\end{equation}
Коллективная точка $c=1$ даёт $0.252005943190\ldots$ при $x=e^{-2}$,
а независимая точка $c=0$ не даёт расщепления вообще.

Даже максимальная вполне положительная точка остаётся ниже сравнительной
мишени $977/3490$. Действительно,
\begin{equation}
 \frac{977}{3490}-v_1(x)
 =\frac{366375x^3-711068x^2+203619x+358774}
 {3490(375x^3+3916x^2+7267x+2782)}.
 \label{eq:v8-baryon-cp-cone-target-gap}
\end{equation}
При $x=e^{-2}<1/7$ числитель строго больше
$358774-711068/49>0$.

\begin{nogocriterion}[Граница отношений $3$ и $3/2$]
Отношения подтверждены точной теорией представлений, но не являются
селектором общего шума. Для получения единственного барионного оператора
нужно независимо вывести ковариацию сред между тремя кварковыми копиями.
Без этого коллективная точка $c=1$ остаётся дополнительным условием.
\end{nogocriterion}

\begin{protocol}
\begin{enumerate}
 \item поле вычислений: \textbf{$\mathbb Q(x,c)$ без чисел с плавающей точкой};
 \item ранг слабых проекторов: \textbf{$4+4$};
 \item спектр $\mathbf I^2$: \textbf{$15/4$ и $3/4$};
 \item аддитивный коэффициент $3$: \textbf{точно};
 \item отношение расщеплений $3/2$: \textbf{точно};
 \item область полной положительности: \textbf{$[-1/2,1]$};
 \item одночастичный генератор на всём отрезке: \textbf{одинаков};
 \item единственность трёхчастичного подъёма: \textbf{запрещена};
 \item коллективный шум $c=1$: \textbf{условие, не вывод};
 \item физическая теорема о массах: \textbf{не получена};
 \item следующий гейт: \textbf{происхождение общей корреляции среды,
 выполнен в следующей главе}.
\end{enumerate}
\end{protocol}

Машинный сертификат сохранён в
\path{s2t_v8_baryon_three_particle_lift_normalization_gate_results.json}.